\section{第 10.1 节：一维空间中的多个粒子}
\label{section:QM-C10-S01}
\studymeta
  {QM-C10-S01}
  {2026-10-09 至 2026-10-10}
  {Shankar, \textit{Principles of Quantum Mechanics}, Second Edition}
  {第 10 章 \S 10.1 N Particles in One Dimension；印刷页 pp.~247--259 开头，PDF pp.~257--269 开头；式 (10.1.1)--(10.1.41)}
  {四块讲解与两道选定习题已核对；笔记已确认}

\textbf{范围与主线。}从两粒子的联合波函数出发，构造 tensor product（张量积，教材也称 direct product），再讨论无相互作用时的演化，以及相互作用二体问题的质心分离。最后推广到 $N$ 个一维粒子，不进入 \S 10.2。本节无独立编号 Example；选做 Exercise 10.1.2 第 (3) 问的一种方法，以及 10.1.3 第 (1) 问，完整解答见课程题目板块。以下先按粒子编号构造空间，不处理全同粒子的交换对称性约束。

\notetopic{QM-C10-S01-T01}{联合位置基、联合密度与边缘密度}
\textbf{物理空间与构型空间。}两个粒子都在同一条直线上运动，下标 $1,2$ 是粒子编号，而非空间方向。描述整个系统的位置需要 $(x_1,x_2)$，因此 configuration space（构型空间）为二维；物理空间仍是一维，相应波函数空间 $L^2(\real^2)$ 则是无限维的 Hilbert 空间。

\textbf{正则关系与共同本征基。}推广单粒子正则关系，
\begin{equation}
  [X_i,P_j]=\ii\hbar\delta_{ij}I,\qquad
  [X_i,X_j]=[P_i,P_j]=0,\qquad i,j=1,2.
\end{equation}
$X_1,X_2$ 可同时对角化，联合位置基满足
\begin{equation}
  X_1\ket{x_1,x_2}=x_1\ket{x_1,x_2},\qquad
  X_2\ket{x_1,x_2}=x_2\ket{x_1,x_2}.
\end{equation}
$\ket{x_1,x_2}$ 表示两个位置条件\textbf{同时成立}，不是一个粒子在两个位置之间的叠加。其广义正交归一与完备关系为
\begin{align}
  \braket{x_1',x_2'}{x_1,x_2}
    &=\delta(x_1'-x_1)\delta(x_2'-x_2),\\
  \int_{\real^2}\ketbra{x_1,x_2}{x_1,x_2}\dd x_1\dd x_2&=I.
\end{align}
不同粒子的基本算符对易不等于它们没有相互作用；是否相互作用要看 Hamiltonian。

\textbf{联合波函数。}整个系统仍由一个态矢量描述，位置表象为
\begin{equation}
  \Psi(x_1,x_2,t)=\braket{x_1,x_2}{\Psi(t)},\qquad
  \ket{\Psi(t)}=\int_{\real^2}\Psi(x_1,x_2,t)\ket{x_1,x_2}\dd x_1\dd x_2.
\end{equation}
算符在此表象中的作用是
\begin{equation}
  X_i\Psi=x_i\Psi,\qquad P_i\Psi=-\ii\hbar\frac{\partial\Psi}{\partial x_i}.
\end{equation}
偏导数只对相应粒子的位置变量作用，求导时固定另一个变量。

\textbf{概率解释与归一化。}定义 joint probability density（联合概率密度）
\begin{equation}
  \rho(x_1,x_2,t)=|\Psi(x_1,x_2,t)|^2,\qquad
  \int_{\real^2}\rho(x_1,x_2,t)\dd x_1\dd x_2=1.
\end{equation}
同时在区间 $A$ 内找到粒子 1、在区间 $B$ 内找到粒子 2 的概率为
\begin{equation}
  \Pr(x_1\in A,\,x_2\in B)
  =\int_A\dd x_1\int_B\dd x_2\,\rho(x_1,x_2,t).
\end{equation}
\textbf{对话补充：边缘密度。}不限制粒子 2 的位置时，对它所有可能的位置积分：
\begin{equation}
  \boxed{\rho_1(x_1,t)=\int_{\real}|\Psi(x_1,x_2,t)|^2\dd x_2}.
  \label{eq:QM-C10-S01-marginal}
\end{equation}
这正是经典联合密度的 marginalization（边缘化）。相加的是概率，不是先积分振幅再取模平方；$|\Psi(x_1,0,t)|^2$ 只是一条截面，既不是边缘密度，也不是已归一化的条件密度。

\notetopic{QM-C10-S01-T02}{张量积空间、乘积态与单粒子算符}
\textbf{张量积要构造什么。}内积给出标量，外积 $\ketbra{u}{v}$ 给出算符；张量积 $\ket u\otimes\ket v$ 则把两个空间的态组合为\textbf{联合空间中的态}：
\begin{equation}
  \hilbert=\hilbert_1\otimes\hilbert_2,\qquad
  \ket{x_1,x_2}=\ket{x_1}\otimes\ket{x_2}.
\end{equation}
张量积对两个因子分别线性，例如
\begin{equation}
  (a\ket u+b\ket{u'})\otimes\ket v
   =a\ket u\otimes\ket v+b\ket{u'}\otimes\ket v.
\end{equation}
乘积向量的内积定义为
\begin{equation}
  (\bra{u'}\otimes\bra{v'})(\ket u\otimes\ket v)
    =\braket{u'}u\,\braket{v'}v.
\end{equation}
因此归一化向量的乘积仍归一化；两组正交归一基 $\{\ket{u_i}\}$、$\{\ket{v_j}\}$ 组成正交归一乘积基 $\{\ket{u_i}\otimes\ket{v_j}\}$。有限维时，每种配对都是一个基向量，故
\begin{equation}
  \boxed{\dim(\hilbert_1\otimes\hilbert_2)=d_1d_2}.
\end{equation}
这不是维数为 $d_1+d_2$ 的 direct sum（直和）。无限维时，Hilbert 张量积包含有限乘积线性组合的范数极限。

\textbf{乘积基完备，不等于所有态都是单个乘积。}以离散正交归一基书写，任意态可展开为
\begin{equation}
  \ket\Psi=\sum_{i,j}C_{ij}\ket{u_i}\otimes\ket{v_j},\qquad
  \sum_{i,j}|C_{ij}|^2=1.
  \label{eq:QM-C10-S01-product-basis}
\end{equation}
对应位置波函数为 $\Psi(x_1,x_2)=\sum_{i,j}C_{ij}u_i(x_1)v_j(x_2)$。教材的坐标论证是：先固定 $x_1$，把关于 $x_2$ 的函数展开为 $\sum_j a_j(x_1)v_j(x_2)$，再将每个 $a_j(x_1)$ 展开为 $\sum_i C_{ij}u_i(x_1)$；所得就是双重展开。对一般 $L^2$ 函数，此处按均方收敛理解。

若联合态是 product state（乘积态），则
\begin{equation}
  \left(\sum_i a_i\ket{u_i}\right)\otimes
  \left(\sum_j b_j\ket{v_j}\right)
    =\sum_{i,j}a_i b_j\ket{u_i}\otimes\ket{v_j}.
\end{equation}
故能分解成单个乘积，当且仅当 $C_{ij}=a_i b_j$；一般系数不满足此条件。对\textbf{纯态}，不可如此分解称为 entangled state（纠缠态）。这里仅补充术语，不展开混合态的可分性判据。

\textbf{对话补充：位置概率独立还不够。}归一化乘积波函数 $u(x_1)v(x_2)$ 必然给出两个位置测量结果的独立分布，但反过来不成立，因为模平方丢失相位。取两个处处非零的归一化高斯函数及实数 $\gamma\ne0$，$\gamma$ 的量纲为长度的负二次方，令
\begin{equation}
  \Psi(x_1,x_2)=u(x_1)v(x_2)\ee^{\ii\gamma x_1x_2}.
\end{equation}
它仍归一化，且 $|\Psi|^2=|u|^2|v|^2$，但对不同的实数 $a,b$，
\begin{equation}
  \frac{\Psi(x_1,a)}{\Psi(x_1,b)}
   =\frac{v(a)}{v(b)}\ee^{\ii\gamma x_1(a-b)}.
\end{equation}
该比值依赖 $x_1$；若 $\Psi=f(x_1)g(x_2)$，比值应为与 $x_1$ 无关的 $g(a)/g(b)$，矛盾。因此\textbf{可分性是完整复振幅的性质，不能仅由一种测量基下的概率独立来判断}。

\textbf{算符的提升与乘法。}若 $A$ 作用于 $\hilbert_1$、$B$ 作用于 $\hilbert_2$，定义
\begin{equation}
  (A\otimes B)(\ket u\otimes\ket v)=(A\ket u)\otimes(B\ket v),
\end{equation}
并线性延拓到一般态。只操作粒子 1 应写 $A\otimes I_2$，只操作粒子 2 应写 $I_1\otimes B$。在乘积基上先后作用可得
\begin{align}
  (A\otimes B)(C\otimes D)&=(AC)\otimes(BD),\\
  [A\otimes I_2,I_1\otimes B]&=0,\\
  [A\otimes I_2,C\otimes I_2]&=[A,C]\otimes I_2.
\end{align}
同一因子内部仍须保留算符顺序；无界算符的上述运算限于所需乘积有定义的共同定义域。教材简写的 $X_1,P_2$，在联合空间中实际是 $X^{(1)}\otimes I_2$、$I_1\otimes P^{(2)}$。

相应矩阵元满足
\begin{equation}
  (\bra{u_i}\otimes\bra{v_j})(A\otimes B)
  (\ket{u_k}\otimes\ket{v_l})=A_{ik}B_{jl}.
\end{equation}
\relatedexercise{QM-C10-S01-E01}{教材习题 10.1.2 的矩阵构造} 展示如何在固定的有序乘积基下使用这一规则。

\notetopic{QM-C10-S01-T03}{联合演化与无相互作用时的分离}
\textbf{一般两粒子方程。}对本节的非相对论模型，
\begin{equation}
  H=\frac{P_1^2}{2m_1}+\frac{P_2^2}{2m_2}+V(X_1,X_2),
  \qquad \ii\hbar\partial_t\ket{\Psi(t)}=H\ket{\Psi(t)}.
\end{equation}
位置表象给出
\begin{equation}
  \ii\hbar\frac{\partial\Psi}{\partial t}
  =\left[-\frac{\hbar^2}{2m_1}\frac{\partial^2}{\partial x_1^2}
         -\frac{\hbar^2}{2m_2}\frac{\partial^2}{\partial x_2^2}
         +V(x_1,x_2)\right]\Psi.
\end{equation}
波函数的变量增加了，Schr\"odinger 演化规律本身未改变。以下讨论不显含时间、具有适当自伴定义域的 Hamiltonian。

\textbf{无相互作用的条件。}若 $V(x_1,x_2)=V_1(x_1)+V_2(x_2)$，则
\begin{equation}
  \boxed{H=H_1\otimes I_2+I_1\otimes H_2},\qquad
  H_i=\frac{(P^{(i)})^2}{2m_i}+V_i(X^{(i)}).
  \label{eq:QM-C10-S01-separable-h}
\end{equation}
允许各粒子受到单独的外势，但没有彼此间的相互作用；这里动能已按粒子分开。设
\begin{equation}
  H_1\ket{u_n}=\epsilon_n\ket{u_n},\qquad
  H_2\ket{v_m}=\eta_m\ket{v_m}.
\end{equation}
利用张量积算符作用，
\begin{align}
  H(\ket{u_n}\otimes\ket{v_m})
  &=(H_1\ket{u_n})\otimes\ket{v_m}
     +\ket{u_n}\otimes(H_2\ket{v_m})\notag\\
  &=(\epsilon_n+\eta_m)(\ket{u_n}\otimes\ket{v_m}).
\end{align}
所以可选取乘积能量本征基，且 $E_{nm}=\epsilon_n+\eta_m$。\textbf{若不同配对给出同一总能量，它们的线性组合仍是能量本征态，却未必是乘积态。}不能把“一组本征基可取乘积形式”误读成“所有本征态都必须是单个乘积”。

\textbf{坐标形式就是分离变量。}尝试 $\Psi_E(x_1,x_2)=u(x_1)v(x_2)$，则
\begin{equation}
  (H_1u)v+u(H_2v)=Euv
  \quad\Longrightarrow\quad
  \frac{H_1u}{u}+\frac{H_2v}{v}=E.
\end{equation}
相除只在 $u,v\ne0$ 的区域进行。第一项仅依赖 $x_1$，第二项仅依赖 $x_2$；分别固定另一变量，便知二者必须各为常数 $\epsilon,\eta$，满足 $\epsilon+\eta=E$。于是恢复两个单粒子本征方程，节点处回到未相除的方程处理。张量积证明不需要除以波函数。

\textbf{传播算符与一般初态。}两个提升后的单粒子 Hamiltonian 对易，因此
\begin{equation}
  \boxed{U(t)=\ee^{-\ii Ht/\hbar}
   =\ee^{-\ii H_1t/\hbar}\otimes\ee^{-\ii H_2t/\hbar}}.
\end{equation}
在离散谱记号下，一般初态和演化为
\begin{align}
  \ket{\Psi(0)}&=\sum_{n,m}C_{nm}\ket{u_n}\otimes\ket{v_m},\\
  \ket{\Psi(t)}&=\sum_{n,m}C_{nm}
       \ee^{-\ii(\epsilon_n+\eta_m)t/\hbar}
       \ket{u_n}\otimes\ket{v_m}.
\end{align}
连续谱部分改成积分。初始乘积态保持为两个单粒子演化态的乘积；初始纠缠态不会被这种局域酉演化变成乘积态，否则施加同样可分解的逆演化便会推出初态也是乘积态。

\textbf{两种可分离性。}Hamiltonian 可分离描述动力学的结构；给定纯态可分离描述其复振幅能否分解。二者不等价，也不能由不同粒子基本算符对易直接推断。

\notetopic{QM-C10-S01-T04}{质心、相对运动与多个自由度}
\textbf{相互作用二体问题的特殊结构。}一般 $V(x_1,x_2)$ 不能分离，但若没有外势且 $V=V(x_1-x_2)$，则可使用第 2.3 节的变量。令
\begin{equation}
  M=m_1+m_2,\qquad \mu=\frac{m_1m_2}{M},
\end{equation}
并定义
\begin{align}
  X_{\mathrm{CM}}&=\frac{m_1X_1+m_2X_2}{M},
  &X_r&=X_1-X_2,\\
  P_{\mathrm{CM}}&=P_1+P_2,
  &P_r&=\frac{m_2P_1-m_1P_2}{M}.
\end{align}
$P_r$ 的经典对应物为 $\mu(\dot x_1-\dot x_2)$，不是一般的 $p_1-p_2$；$\mu$ 是描述相对运动的约化质量，不是另一个真实粒子的质量。

\textbf{验证正则关系。}利用原来的对易关系，
\begin{align}
  [X_{\mathrm{CM}},P_{\mathrm{CM}}]
    &=\frac{m_1+m_2}{M}\ii\hbar I=\ii\hbar I,\\
  [X_r,P_r]&=\frac{m_2+m_1}{M}\ii\hbar I=\ii\hbar I,\\
  [X_{\mathrm{CM}},P_r]
    &=\frac{m_1m_2-m_2m_1}{M^2}\ii\hbar I=0,\\
  [X_r,P_{\mathrm{CM}}]&=\ii\hbar I-\ii\hbar I=0.
\end{align}
新位置之间、新动量之间也对易，故得到两套正则自由度。

\textbf{动能中的交叉项抵消。}反解为
\begin{equation}
  P_1=\frac{m_1}{M}P_{\mathrm{CM}}+P_r,\qquad
  P_2=\frac{m_2}{M}P_{\mathrm{CM}}-P_r.
\end{equation}
代入动能，并利用 $[P_{\mathrm{CM}},P_r]=0$，
\begin{align}
  \frac{P_1^2}{2m_1}+\frac{P_2^2}{2m_2}
  &=\frac{m_1+m_2}{2M^2}P_{\mathrm{CM}}^2
    +\left(\frac1M-\frac1M\right)P_{\mathrm{CM}}P_r\notag\\
  &\quad+\frac12\left(\frac1{m_1}+\frac1{m_2}\right)P_r^2
   =\frac{P_{\mathrm{CM}}^2}{2M}+\frac{P_r^2}{2\mu}.
\end{align}
于是
\begin{equation}
  \boxed{H=\underbrace{\frac{P_{\mathrm{CM}}^2}{2M}}_{H_{\mathrm{CM}}}
    +\underbrace{\left[\frac{P_r^2}{2\mu}+V(X_r)\right]}_{H_r}}.
  \label{eq:QM-C10-S01-cm-separation}
\end{equation}
分离的是\textbf{质心与相对运动}，不是原来的两个粒子；相互作用完整地保留在 $V(X_r)$ 中。

\textbf{乘积本征解。}用 $R,r$ 表示新位置变量。一组可选的广义能量本征函数及能量为
\begin{equation}
  \boxed{\Psi_{p_{\mathrm{CM}},\epsilon}(R,r)
   =\frac{\ee^{\ii p_{\mathrm{CM}}R/\hbar}}{\sqrt{2\pi\hbar}}
     \phi_\epsilon(r)},\qquad
  E=\frac{p_{\mathrm{CM}}^2}{2M}+\epsilon,
\end{equation}
其中
\begin{equation}
  \left[-\frac{\hbar^2}{2\mu}\frac{\dd^2}{\dd r^2}+V(r)\right]
       \phi_\epsilon(r)=\epsilon\phi_\epsilon(r).
\end{equation}
质心部分是自由粒子平面波，采用动量 delta 归一化；正规化的联合波包还需适当叠加。将乘积解换回 $x_1,x_2$ 后，一般不能写成 $u(x_1)v(x_2)$，因为两种新变量都混合了原来的粒子坐标。

\textbf{坐标变换不等于更换参考系（对话补充）。}$(x_1,x_2)\mapsto(R,r)$ 是构型空间中的变量变换，未改变物理系统，也未消除相互作用。另选质心参考系才是参考系变换。教材对确定总动量的本征解取 $p_{\mathrm{CM}}=0$，从而专注相对运动；对总动量有分布的波包，惯性系变换可使平均总动量为零，却不能消除它的动量不确定度。

\textbf{推广到 $N$ 个一维粒子。}联合空间为 $\hilbert_1\otimes\cdots\otimes\hilbert_N$，波函数为 $\Psi(x_1,\ldots,x_N)$，归一化使用 $N$ 重位置积分。无相互作用时，乘积本征基与能量相加继续成立。有相互作用时，即使平移不变性允许分离质心，剩余相对自由度通常仍耦合，不能保证全部化成独立单粒子问题。稳定的二次型耦合振子是重要特例，可用 normal coordinates（正常坐标）进一步解耦。

\relatedexercise{QM-C10-S01-E02}{教材习题 10.1.3 的正常模量子化} 将两个耦合粒子改写为频率不同的独立模式。该线性正则变换可在量子化前进行，也可在原坐标量子化后改写微分方程；本次只选做前一路线。对曲线坐标或一般非线性正则变换，不能据此保证直接替换给出唯一量子化；教材提醒先在笛卡尔坐标量子化，再处理坐标变换。

\textbf{一分钟回顾。}联合波函数描述整个系统；边缘密度对未观察的位置积分。张量积的维数相乘，乘积基完备不等于任意态都是乘积；相位可以使位置分布独立的纯态仍然纠缠。无相互作用时总能量相加、传播算符分解，但初态未必可分。相互作用二体问题有时能按质心与相对运动分离，说明可分离性必须注明相对于哪组自由度。
